\documentclass{article}

\usepackage{jrm-roteiro}

\title{Roteiro da aula prática 06\\Pedidos}
\author{BCC222 --- Programação funcional}
\date{}

\begin{document}

\maketitle

\begin{abstract}
Esta aula prática propõe a implementação de funções para um pequeno sistema
de preparação de pedidos de uma loja virtual. A atividade integra tuplas,
listas, \haskellinline{Maybe}, listas de associação e sinônimos de tipo em um
mesmo contexto, permitindo que os conceitos do capítulo 6 sejam exercitados
de forma aplicada.
\end{abstract}

\section{Objetivos}

Ao final desta atividade, espera-se que você seja capaz de:

\begin{itemize}
  \item manipular dados organizados em tuplas e listas;
  \item usar \haskellinline{Maybe} para tratar a possível ausência de valor;
  \item consultar informações em listas de associação usando
  \haskellinline{lookup};
  \item empregar sinônimos de tipo para tornar a modelagem mais legível;
  \item aplicar a ideia de recursão sobre listas em um exercício introdutório.
\end{itemize}

\section{Contexto}

Uma loja virtual precisa representar pedidos de clientes, consultar
informações básicas sobre cada pedido e tratar situações em que o valor total
ainda não foi calculado. Além disso, a loja mantém uma tabela de descontos
por cliente.

Nesta atividade, o contexto de \emph{pedidos} será usado para integrar, em um
único problema, os conceitos estudados no capítulo 6.

\section{Modelagem fornecida}

A atividade utiliza os seguintes tipos:

\begin{haskellcode}
type Cliente = String
type NomeProduto = String
type Quantidade = Int
type ValorTotal = Double

type ItemPedido = (NomeProduto, Quantidade)
type Pedido = (Cliente, [ItemPedido], Maybe ValorTotal)
type TabelaDescontos = [(Cliente, Double)]
\end{haskellcode}

Para acessar os componentes de uma tripla, você pode usar as funções
\haskellinline{fst3}, \haskellinline{snd3} e \haskellinline{thd3}, do módulo
\haskellinline{Data.Tuple.Extra}.

\section{Tarefas}

Implemente, no arquivo \texttt{src/Pedidos.hs}, as seguintes funções.

\subsection*{1. Acesso aos componentes do pedido}

\begin{haskellcode}
clientePedido :: Pedido -> Cliente
itensPedido :: Pedido -> [ItemPedido]
totalPedido :: Pedido -> Maybe ValorTotal
\end{haskellcode}

\subsection*{2. Informações derivadas do pedido}

\begin{haskellcode}
quantidadeTiposItens :: Pedido -> Int
temTotalCalculado :: Pedido -> Bool
totalOuPadrao :: ValorTotal -> Pedido -> ValorTotal
\end{haskellcode}

\subsection*{3. Consulta em tabela de descontos}

\begin{haskellcode}
consultarDesconto :: Cliente -> TabelaDescontos -> Double
clienteTemDesconto :: Cliente -> TabelaDescontos -> Bool
\end{haskellcode}

\subsection*{4. Recursão sobre a lista de itens}

\begin{haskellcode}
quantidadeTotalUnidades :: [ItemPedido] -> Int
\end{haskellcode}

Nesta última função, a sugestão é usar recursão. Pense no caso da lista vazia
e no caso em que há um primeiro item e o restante da lista.

\section{Orientações}

\begin{itemize}
  \item Edite apenas o arquivo \texttt{src/Pedidos.hs}.
  \item Use o \texttt{README.md} do repositório como referência principal do
  enunciado.
  \item O fluxo de trabalho com GitHub Classroom, Git, Cabal e GHCi não será
  reexplicado aqui; consulte o roteiro da atividade prática do capítulo 5, se
  necessário.
  \item Para abrir o ambiente interativo, use \texttt{cabal repl}.
  \item Para executar os testes, use \texttt{cabal test}.
\end{itemize}

\section{Avaliação}

A correção automática está dividida por função, permitindo avaliação
granular. A distribuição de pontos é:

\begin{itemize}
  \item \haskellinline{clientePedido}: 1 ponto
  \item \haskellinline{itensPedido}: 1 ponto
  \item \haskellinline{totalPedido}: 1 ponto
  \item \haskellinline{quantidadeTiposItens}: 1 ponto
  \item \haskellinline{temTotalCalculado}: 1 ponto
  \item \haskellinline{totalOuPadrao}: 1 ponto
  \item \haskellinline{consultarDesconto}: 1 ponto
  \item \haskellinline{clienteTemDesconto}: 1 ponto
  \item \haskellinline{quantidadeTotalUnidades}: 2 pontos
\end{itemize}

\section{Entrega}

\begin{itemize}
  \item \textbf{Data regular de entrega:} até 28 de abril.
  \item \textbf{Data limite com penalização:} até 5 de maio.
  \item Após a data regular, haverá 9 dias de tolerância, com penalização de
  10\% por dia de atraso.
\end{itemize}

\section{Observação final}

Embora a maior parte da atividade possa ser resolvida com funções já
estudadas, a última função serve como uma primeira aproximação à ideia de
recursão sobre listas, que será estudada formalmente em um capítulo futuro.

\end{document}
